%; whizzy paragraph -pdf xpdf -latex ./whizzypdfptex.sh
%; whizzy-paragraph "^\\\\begin{frame}\\|\\\\emtext"
% latex beamer presentation.
% platex, latex-beamer でコンパイルすることを想定。 

%     Tokyo Debian Meeting resources
%     Copyright (C) 2012 Junichi Uekawa
%     Copyright (C) 2016 Nobuhiro Iwamatsu

%     This program is free software; you can redistribute it and/or modify
%     it under the terms of the GNU General Public License as published by
%     the Free Software Foundation; either version 2 of the License, or
%     (at your option) any later version.

%     This program is distributed in the hope that it will be useful,
%     but WITHOUT ANY WARRANTY; without even the implied warreanty of
%     MERCHANTABILITY or FITNESS FOR A PARTICULAR PURPOSE.  See the
%     GNU General Public License for more details.

%     You should have received a copy of the GNU General Public License
%     along with this program; if not, write to the Free Software
%     Foundation, Inc., 51 Franklin St, Fifth Floor, Boston, MA  02110-1301 USA

\documentclass[cjk,dvipdfmx,12pt]{beamer}
\usetheme{Tokyo}
\usepackage{monthlypresentation}

%  preview (shell-command (concat "evince " (replace-regexp-in-string "tex$" "pdf"(buffer-file-name)) "&")) 
%  presentation (shell-command (concat "xpdf -fullscreen " (replace-regexp-in-string "tex$" "pdf"(buffer-file-name)) "&"))
%  presentation (shell-command (concat "evince " (replace-regexp-in-string "tex$" "pdf"(buffer-file-name)) "&"))

%http://www.naney.org/diki/dk/hyperref.html
%日本語EUC系環境の時
\AtBeginDvi{\special{pdf:tounicode EUC-UCS2}}
%シフトJIS系環境の時
%\AtBeginDvi{\special{pdf:tounicode 90ms-RKSJ-UCS2}}

\newenvironment{commandlinesmall}%
{\VerbatimEnvironment
  \begin{Sbox}\begin{minipage}{1.0\hsize}\begin{fontsize}{8}{8} \begin{BVerbatim}}%
{\end{BVerbatim}\end{fontsize}\end{minipage}\end{Sbox}
  \setlength{\fboxsep}{8pt}
% start on a new paragraph

\vspace{6pt}% skip before
\fcolorbox{dancerdarkblue}{dancerlightblue}{\TheSbox}

\vspace{6pt}% skip after
}
%end of commandlinesmall

\title{apacheにdjangoのwebアプリを\\デプロイしてみよう}
\subtitle{OSC 2018 Tokyo/Spring\\東京エリアDebian勉強会/Debian JP Project\\第160回出張勉強会}
\author{杉本 典充\\ dictoss@live.jp}
\date{2018年02月24日}
\logo{\includegraphics[width=8cm]{image200607/openlogo-light.eps}}

\begin{document}

\begin{frame}
\titlepage{}
\end{frame}

\begin{frame}{Agenda}
  \begin{itemize}
  \item Debian と Debian JPについて
  \item Debian と webサーバ
  \item DebianでDjangoなwebアプリを動かす
    \begin{itemize}
    \item 実行環境
    \item webサーバ、WSGIサーバのインストール
    \item webアプリケーションの配置
    \item mod\_wsgiの設定
    \item apache2の設定
    \item 動作確認
    \end{itemize}
  \item まとめ
  \item 今後のイベントと情報入手先
  \end{itemize}
\end{frame}

%-------------------

\section{Debian と Debian JPについて}

\begin{frame}
  \begin{center}\Huge{Debian と Debian JPについて}\end{center}
\end{frame}

\subsection{Debian とは？}

\begin{frame}{Debian とは？}
  
\begin{itemize}
\item \url{https://www.debian.org/}
\item Debianはフリー/オープンなオペレーティングシステム (OS)を作成しようとするボランティアベースのプロジェクト。
\item 自分たちの考えるフリーという言葉に関する定義、開発目的、パッケージングポリシーを厳格に決めている。
\item 世界中に1000人以上の開発者がおり、他のディストリビューションのベースとして採用。
\item 約2年毎にリリースが行われ、多くのパッケージとアーキテクチャをサポート。
\item 年に一度、Debconfという開発者が集うイベントがある。Debconf18は7/21-7/27に台湾で開催予定。
\item サーバ、PC、組み込みなど様々なところで利用されているLinuxディストリビューション。
\end{itemize}

\end{frame}


\subsection{Debian JP Project と Debian勉強会}

\begin{frame}{Debian JP Project と Debian勉強会}
  \begin{itemize}
  \item Debian JP Project
    \begin{itemize}
    \item \url{https://www.debian.or.jp/}
    \item 日本でDebianを普及させることを目的とした任意団体。
    \item Debian の日本語による情報発信、ユーザとの情報交換、Debian 開発者、パッケージメンテナの育成など。
    \end{itemize}
  \item Debian勉強会
    \begin{itemize}
    \item 2005年1月開始。
    \item 東京と関西で毎月開催しているDebian開発者、ユーザによる勉強会。
    \item \url{https://tokyodebian.alioth.debian.org/}
    \item \url{https://wiki.debian.org/KansaiDebianMeeting}
    \item 東京エリアは第3土曜日、関西エリアは第4日曜日に開催することが多い。
    \end{itemize}
  \end{itemize}
\end{frame}

%-----------------------

\section{Debianとwebサーバ}

\begin{frame}
  \begin{center}\Huge{Debianとwebサーバ}\end{center}
\end{frame}


\subsection{Debianのwebサーバで言われること}

\begin{frame}{Debianのwebサーバで言われること}% [containsverbatim]
  
\begin{itemize}
\item webサーバ用途でDebianはあまり使われていないのではないか
\item webサーバのパッケージは多数あり、どれをインストールするとよいのかわからない
\item webサーバの設定ファイルのディレクトリ構造がRHELやtarballでインストールしたときと異なるためよくわからない
\end{itemize}

そんな疑問にお答えするため、今回はdebian 9 の apache2 及び pythonのWebアプリケーションサーバ"WSGI"の設定についてお話しします。

\end{frame}


\subsection{webサーバになぜDebianを選択するのか}

\begin{frame}{webサーバになぜDebianを選択するのか}% [containsverbatim]
\begin{itemize}
  \item Linuxのwebサーバの稼働OSシェア\footnote{\url{https://w3techs.com/technologies/details/os-linux/all/all}}
  \begin{itemize}
    \item 1位 Ubuntu 39.1 \%
    \item 2位 Debian 30.8 \%
    \item 3位 CentOS 20.6 \%
    \item 4位 Red Hat 2.9 \%
  \end{itemize}

  \item 月100万アクセス以上のwebサイトのwebサーバシェア\footnote{\url{https://news.netcraft.com/archives/2018/01/19/january-2018-web-server-survey.html}}
  \begin{itemize}
    \item 1位 Apache 37.03 \% %% 2012年くらいから下がり始めた
    \item 2位 nginx 30.46 \% %% apacheの減少がnginxに移っている
    \item 3位 Microsoft IIS Web Server 9.57 \%      
  \end{itemize}
\end{itemize}

\end{frame}


\subsection{Debianのwebサーバのパッケージ}

\begin{frame}{Debianのwebサーバのパッケージ}% [containsverbatim]

\begin{itemize}
  \item apache2 (apache-2.4.25)
  \item nginx (nginx-1.10.3、stretch-backports\footnote{backportsとは、安定版に比較的新しいソフトウェアを提供する仕組み}はnginx-1.13.3)
  \item lighttpd (lighttpd-1.4.45)
\end{itemize}

\end{frame}


\subsection{DebainのWSGIサーバのパッケージ}

\begin{frame}{DebainのWSGIサーバのパッケージ}% [containsverbatim]
  
\begin{itemize}
\item WSGI = Web Server Gateway Interface
\item PEP 333にて定義\footnote{\url{https://www.python.org/dev/peps/pep-0333/}}
\item WSGIを実装したwebアプリケーションサーバ
  \begin{itemize}
  \item libapache2-mod-wsgi、libapache2-mod-wsgi-py3 パッケージ
    \begin{itemize}
    \item apache2の拡張モジュールmod\_wsgi
    \item apache2のみ利用可能
    \end{itemize}
  \item uwsgi パッケージ
    \begin{itemize}
    \item nginx、apache2で利用可能
    \end{itemize}
  \item python-gunicorn、python3-gunicorn パッケージ
    \begin{itemize}
    \item nginx、apache2で利用可能
    \end{itemize}
  \end{itemize}
\item 性能ベンチマークの例（2016年5月）\footnote{\url{https://blog.appdynamics.com/engineering/a-performance-analysis-of-python-wsgi-servers-part-2/}}
\end{itemize}

\end{frame}


\section{DebianでDjangoなwebアプリを動かす}

\begin{frame}
  \begin{center}\Huge{DebianでDjangoなwebアプリを動かす}\end{center}
\end{frame}


\subsection{実行環境}

\begin{frame}{実行環境}% [containsverbatim]

\begin{itemize}
\item 作成するアプリケーションの実行環境
  \begin{itemize}
    \item Debian GNU/Linux 9 Stretch amd64
    \item apache-2.4
    \item mod\_wsgi-4.5
    \item python-3.5
    \item sqlite3
    \item django-2.0 (debパッケージを利用せずpipでインストール)
  \end{itemize}
\item 作成する実行環境の前提
  \begin{itemize}
  \item virtualenvは使わず、システムワイドにpythonライブラリをインストールする
  \end{itemize}
\item 情報
  \begin{itemize}
  \item \url{https://docs.djangoproject.com/ja/2.0/howto/deployment/wsgi/modwsgi/}
  \end{itemize}
\end{itemize}

\end{frame}


\subsection{webアプリケーションの開発}

\begin{frame}[containsverbatim]{webアプリケーションの開発}

\begin{itemize}
\item Djangoのアプリケーション開発のチュートリアル
  \begin{itemize}
  \item \url{https://docs.djangoproject.com/ja/2.0/intro/}
  \end{itemize}
\item サンプルソースコード
  \begin{itemize}
  \item \url{https://github.com/dictoss/django-tutorial}
  \end{itemize}
\end{itemize}

\end{frame}


\subsection{webサーバ、WSGIサーバのインストール}

\begin{frame}[containsverbatim]{webサーバ、WSGIサーバのインストール}
 
\begin{commandlinesmall}
# apt update
# apt install python3.5 python3-pip
# apt install apache2 libapache2-mod-wsgi-py3
# apt install sqlite3
\end{commandlinesmall}

\begin{commandlinesmall}
# pip3 install django==2.0.2
\end{commandlinesmall}

\begin{itemize}
\item 注意点
  \begin{itemize}
  \item libapache2-mod-wsgi-py3とlibapache2-mod-wsgiはパッケージの制約で片方しかインストールできない
  \item つまり、mod\_wsgiではpython2.7、python3.5のアプリを同居できない
  \end{itemize}
\end{itemize}

\end{frame}


\subsection{webアプリケーションの配置}

\begin{frame}[containsverbatim]{webアプリケーションの配置}

\begin{commandlinesmall}
# apt install git
\end{commandlinesmall}

\begin{commandlinesmall}
$ git clone https://github.com/dictoss/django-tutorial.git
$ cd django-tutorial/2.0
$ su
# mkdir -p /var/www/django-tutorial
# cp -r mysite /var/www/django-tutorial/
# chown -fR www-data:www-data /var/www/django-tutorial
\end{commandlinesmall}

\end{frame}


\subsection{mod\_wsgiの設定}

\begin{frame}[containsverbatim]{mod\_wsgiの設定（1）}

mod\_wsgi全体の動作設定

\begin{commandline}
# vi /etc/apache2/mods-available/wsgi.conf
\end{commandline}

特に変更しなくても動きます。
詳しい設定はこちら。\footnote{\url{http://modwsgi.readthedocs.io/}}
\end{frame}


\begin{frame}[containsverbatim]{mod\_wsgiの設定（2）}

mod\_wsgiにおけるアプリケーションの動作モードは2種類\footnote{\url{http://modwsgi.readthedocs.io/en/latest/user-guides/reloading-source-code.html}}
  
\begin{itemize}
  \item Daemon Mode（デフォルト）
  \item Embedded Mode（WSGIRestrictEmbedded Onを指定）
\end{itemize}

\begin{table}[htb]
  \caption{mod\_wsgiの動作モードの比較}
  \begin{tabular}{|l|c|c|}
    \hline
    ― & Embedded Mode & Daemon Mode \\
    \hline \hline
    メモリ使用量 & 大 & 中 \\
    \hline
    CPU使用量 & 大 & 中 \\
    \hline
    性能 & 高 & 中 \\
    \hline
    アプリ再読み込み & restart apache2 & \shortstack{reload apache2 \\ or touch wsgi.py} \\
    \hline
    apache MPM & prefork & \shortstack{prefork \\ worker \\ event} \\
    \hline
  \end{tabular}
\end{table}

\end{frame}


\begin{frame}[containsverbatim]{mod\_wsgiの設定（3）}

アプリケーション固有の設定（Daemon Mode利用時の設定）

\begin{commandlinesmall}
# vi /etc/apache2/conf-available/wsgi_mysite.conf

Alias /static/polls ``/var/www/django-tutorial/mysite/polls/static/polls''
<Directory ``/var/www/django-tutorial/mysite/polls/static/polls''>
  Require all granted
</Directory>

Alias /static/admin \
``/usr/local/lib/python3.5/dist-packages/django/contrib/admin/static/admin''
<Directory \
``/usr/local/lib/python3.5/dist-packages/django/contrib/admin/static/admin''>
  Require all granted
</Directory>
    
WSGIScriptAlias /mysite /var/www/django-tutorial/mysite/mysite/wsgi.py

WSGIDaemonProcess mysite user=www-data group=www-data \
processes=2 threads=64 maximum-requests=256 \
python-path=/var/www/django-tutorial/mysite \
display-name=%{GROUP}

WSGIProcessGroup  mysite
\end{commandlinesmall}

\end{frame}


\begin{frame}[containsverbatim]{mod\_wsgiの設定（4）}

\begin{itemize}
\item もしvirtualenvを使う場合は、WSGIDaemonProcessに"python-home"を指定する
\end{itemize}  

\begin{commandlinesmall}
WSGIDaemonProcess mysite user=www-data group=www-data \
processes=2 threads=64 maximum-requests=256 \
python-path=/var/www/django-tutorial/mysite \
python-home=/path/to/dir \
display-name=%{GROUP}
\end{commandlinesmall}

\begin{itemize}
\item WSGIDaemonProcessの詳細パラメータはドキュメント参照\footnote{\url{http://modwsgi.readthedocs.io/en/latest/configuration-directives/WSGIDaemonProcess.html}}
\end{itemize}
    
\end{frame}


\subsection{apache2の設定}

\begin{frame}[containsverbatim]{apache2の設定（1）}

Debianの"/etc/apache2"ディレクトリの構成

\begin{commandlinesmall}
$ tree /etc/apache2
/etc/apache2
├── apache2.conf
├── conf-available
│     ├── （設定ファイルを保存）
├── conf-enabled
│     ├── （設定ファイルの有効/無効設定を保存）
├── envvars
├── magic
├── mods-available
│     ├── （拡張モジュールの設定ファイルを保存）
├── mods-enabled
│     ├── （拡張モジュールの有効/無効設定を保存）
├── ports.conf
├── sites-available
│     ├── （VirtualHostの設定ファイルを保存）
└── sites-enabled
       ├── （VirtualHostの有効/無効設定を保存）
\end{commandlinesmall}
  
\end{frame}


\begin{frame}[containsverbatim]{apache2の設定（2）}

Debianの"/etc/apache2"ディレクトリの解説

\begin{itemize}
  \item apache2.conf, envvars, magic, ports.conf
    \begin{itemize}
    \item 設定や環境変数を定義するファイル
    \end{itemize}
  \item *-available ディレクトリ
    \begin{itemize}
    \item conf、mods、siteの3種類
    \item 設定ファイルの保存場所
    \end{itemize}
  \item *-enabled ディレクトリ
    \begin{itemize}
    \item conf、mods、siteの3種類
    \item 有効/無効の設定の保存場所
    \item a2enconf と a2disconf でconfの有効/無効を切り替え
    \item a2enmod と a2dismod でmodsの有効/無効を切り替え
    \item a2ensite と a2dissite でsitesの有効/無効を切り替え
    \end{itemize}
\end{itemize}
  
\end{frame}


\begin{frame}[containsverbatim]{apache2の設定（3）}

作成したWSGIアプリケーションの設定を有効化

\begin{commandlinesmall}
# cd /etc/apache2
# a2enconf wsgi_mysite
Enabling conf wsgi_mysite.
To activate the new configuration, you need to run:
systemctl reload apache2  
\end{commandlinesmall}

conf-enabledディレクトリの状態を確認

\begin{commandlinesmall}
# cd /etc/apache2/conf-enabled 
# ls -l wsgi_mysite.conf
lrwxrwxrwx 1 root root 34  2月 18 12:28 wsgi_mysite.conf ->
  ../conf-available/wsgi_mysite.conf
\end{commandlinesmall}

\end{frame}


\subsection{動作確認}

\begin{frame}[containsverbatim]{動作確認}

設定を反映するため、apache2を再起動
  
\begin{commandlinesmall}
# systemctl restart apache2
\end{commandlinesmall}

プロセス確認

\begin{commandlinesmall}
# ps -ef | grep www | grep -v grep
www-data  6507  6506  0  3月18 ?      00:00:00 (wsgi:mysite)     -k start
www-data  6508  6506  0  3月18 ?      00:00:00 (wsgi:mysite)     -k start
www-data  6509  6506  0  3月18 ?      00:00:00 /usr/sbin/apache2 -k start
www-data  6510  6506  0  3月18 ?      00:00:00 /usr/sbin/apache2 -k start  
\end{commandlinesmall}

アクセス確認

\url{http://localhost/mysite/polls/}

\end{frame}

\section{まとめ}

\begin{frame}[containsverbatim]{まとめ}

  \begin{itemize}
  \item Debian系OSはwebサーバとして多く利用されている
  \item webサーバとWSGIサーバのdebパッケージの紹介
  \item apache2＋mod\_wsgiのインストール方法
  \item mod\_wsgiで動作するアプリケーションの設定ファイルの作成方法
  \item debianのapache2のディレクトリ構成の解説
  \end{itemize}

\end{frame}

%----------------

\section{今後のイベントと情報入手先}

\begin{frame}{今後のイベントと情報入手先}

  \begin{itemize}
  \item 第132関西Debian勉強会
    \begin{itemize}
      \item 2018/02/25（日） 13:30 - 17:00
      \item 大阪　福島区民センター
      \item \url{https://wiki.debian.org/KansaiDebianMeeting/20180225}
    \end{itemize}
    
  \item 第161回東京エリアDebian勉強会
    \begin{itemize}
      \item 2018/03/24（土） 14:00 - 18:00
      \item \url{http://tokyodebian.alioth.debian.org/2018-03.html}
    \end{itemize}

  \item Debian Project
    \begin{itemize}
    \item Webサイト \url{https://www.debian.org/}
    \item Debian Wiki \url{https://wiki.debian.org/ja}
    \end{itemize}
        
  \item Debian JP Project
    \begin{itemize}
      \item Webサイト \url{https://www.debian.or.jp/}
      \item Twitter @debianjp, @debian\_jp
      \item 雑誌 Software Design 「Debian Hot Topics」
    \end{itemize}
  \end{itemize}

\end{frame}


\end{document}

;;; Local Variables: ***
;;; outline-regexp: "\\([ 	]*\\\\\\(documentstyle\\|documentclass\\|emtext\\|section\\|begin{frame}\\)\\*?[ 	]*[[{]\\|[]+\\)" ***
;;; End: ***
